% GENERATED FILE. Edit canonical lessons, then run npm run generate:books.
% canonical-source-hash: fnv1a64:8b0db2c194f0d4eb
% canonical-lessons: ES-C549-casado, ES-C549-soltero, ES-C549-casar, ES-C549-marido, ES-C549-brindar

\chapter{Casados}
\label{ch:es-casados}
\hlchaptermodality{549}

\begin{chapteropening}
\textbf{By the end of this chapter:} \emph{I can say whether someone is married or single, who married whom, and who raised a glass.}

Complete the last lesson of chapter 549: I can say whether someone is married or single, who married whom, and who raised a glass.
\end{chapteropening}

\section[casado]{casado — married}

\label{lesson:ES-C549-casado}



\begin{quote}
Before the new one: say the Spanish for the toast, then the Spanish for the honeymoon.
\end{quote}

\subsection*{You'll want to know: casado}
\textbf{casado} — \textquotedblleft{}married\textquotedblright{}.

\begin{quote}
Mi hermano está casado con una francesa.
\end{quote}

\begin{quote}
\emph{My brother is married to a Frenchwoman.}
\end{quote}

\emph{Estar casado con} — to be married to.

\begin{cousinweb}
\emph{casado} is built on the participle of \emph{casar} 'to marry', which comes from \emph{casa} 'house', inherited from Latin \emph{casa} 'a hut' (Lewis \& Short). To marry was to set up house.
\end{cousinweb}

\begin{grammarlens}[title={What you've built}]
One more word for telling, explaining and describing at B1.
\end{grammarlens}

\subsection*{Your turn}
\begin{itemize}
\raggedright
  \item \emph{Say it:} \emph{casado}
  \item \emph{Say it:} \emph{casado}, once more
  \item \emph{Say it:} a sentence of your own with \emph{casado}
  \item \emph{Recall:} say the Spanish for the toast, then the Spanish for the honeymoon, then say \emph{casado} again
\end{itemize}

\begin{tcolorbox}[breakable,colback=teal!4,colframe=teal!35!black,title={Before you move on}]
What does casado mean? (\textquotedblleft{}Married\textquotedblright{}.) Say it once more.
\end{tcolorbox}
\section[soltero]{soltero — single, unmarried}

\label{lesson:ES-C549-soltero}



\begin{quote}
Before the new one: say the Spanish for the honeymoon, then the Spanish for married.
\end{quote}

\subsection*{You'll want to know: soltero}
\textbf{soltero} — \textquotedblleft{}single, unmarried\textquotedblright{}.

\begin{quote}
Mi tío nunca se casó; sigue soltero.
\end{quote}

\begin{quote}
\emph{My uncle never married; he is still single.}
\end{quote}

\emph{Un soltero} is a single man; \emph{la despedida de soltero} is a stag party.

\begin{cousinweb}
\emph{soltero} is inherited from Latin \emph{sōlitārius} 'alone, solitary' (EtymDB), built on \emph{sōlus} 'alone'.
\end{cousinweb}

\begin{grammarlens}[title={What you've built}]
One more word for telling, explaining and describing at B1.
\end{grammarlens}

\subsection*{Your turn}
\begin{itemize}
\raggedright
  \item \emph{Say it:} \emph{soltero}
  \item \emph{Say it:} \emph{soltero}, once more
  \item \emph{Say it:} a sentence of your own with \emph{soltero}
  \item \emph{Recall:} say the Spanish for the honeymoon, then the Spanish for married, then say \emph{soltero} again
\end{itemize}

\begin{tcolorbox}[breakable,colback=teal!4,colframe=teal!35!black,title={Before you move on}]
What does soltero mean? (\textquotedblleft{}Single, unmarried\textquotedblright{}.) Say it once more.
\end{tcolorbox}
\section[casar]{casar — to marry (someone to someone)}

\label{lesson:ES-C549-casar}



\begin{quote}
Before the new one: say the Spanish for married, then the Spanish for single.
\end{quote}

\subsection*{You'll want to know: casar}
\textbf{casar} — \textquotedblleft{}to marry (someone to someone)\textquotedblright{}.

\begin{quote}
Los casó el alcalde.
\end{quote}

\begin{quote}
\emph{The mayor married them.}
\end{quote}

Without \emph{se} it is the person who performs the wedding. \emph{Casarse} is to get married.

\begin{cousinweb}
\emph{casar} is built on \emph{casa} 'house', which is inherited from Latin \emph{casa} 'a hut, a cottage' (Lewis \& Short).
\end{cousinweb}

\begin{grammarlens}[title={What you've built}]
One more word for telling, explaining and describing at B1.
\end{grammarlens}

\subsection*{Your turn}
\begin{itemize}
\raggedright
  \item \emph{Say it:} \emph{casar}
  \item \emph{Say it:} \emph{casar}, once more
  \item \emph{Say it:} a sentence of your own with \emph{casar}
  \item \emph{Recall:} say the Spanish for married, then the Spanish for single, then say \emph{casar} again
\end{itemize}

\begin{tcolorbox}[breakable,colback=teal!4,colframe=teal!35!black,title={Before you move on}]
What does casar mean? (\textquotedblleft{}To marry (someone to someone)\textquotedblright{}.) Say it once more.
\end{tcolorbox}
\section[el marido]{el marido — the husband}

\label{lesson:ES-C549-marido}



\begin{quote}
Before the new one: say the Spanish for single, then the Spanish for to marry.
\end{quote}

\subsection*{You'll want to know: el marido}
\textbf{el marido} — \textquotedblleft{}the husband\textquotedblright{}.

\begin{quote}
Su marido trabaja en un hospital.
\end{quote}

\begin{quote}
\emph{Her husband works in a hospital.}
\end{quote}

\emph{El esposo} means the same, a little more formal.

\begin{cousinweb}
\emph{marido} is inherited from Latin \emph{marītus} 'a husband' (EtymDB).
\end{cousinweb}

\begin{grammarlens}[title={What you've built}]
One more word for telling, explaining and describing at B1.
\end{grammarlens}

\subsection*{Your turn}
\begin{itemize}
\raggedright
  \item \emph{Say it:} \emph{el marido}
  \item \emph{Say it:} \emph{el marido}, once more
  \item \emph{Say it:} a sentence of your own with \emph{el marido}
  \item \emph{Recall:} say the Spanish for single, then the Spanish for to marry, then say \emph{el marido} again
\end{itemize}

\begin{tcolorbox}[breakable,colback=teal!4,colframe=teal!35!black,title={Before you move on}]
What does el marido mean? (\textquotedblleft{}The husband\textquotedblright{}.) Say it once more.
\end{tcolorbox}
\section[brindar]{brindar — to toast, to drink to; to offer}

\label{lesson:ES-C549-brindar}



\begin{quote}
Before the new one: say the Spanish for to marry, then the Spanish for the husband.
\end{quote}

\subsection*{You'll want to know: brindar}
\textbf{brindar} — \textquotedblleft{}to toast, to drink to; to offer\textquotedblright{}.

\begin{quote}
Brindamos por el nuevo año.
\end{quote}

\begin{quote}
\emph{We drank a toast to the new year.}
\end{quote}

\emph{Brindar por} — to drink to. \emph{Brindar ayuda} — to offer help.

\begin{cousinweb}
\emph{brindar} is built on \emph{brindis} 'a toast' (EtymDB), which comes from the German words \emph{bring dir's} 'I bring it to you'.
\end{cousinweb}

\begin{grammarlens}[title={What you've built}]
One more word for telling, explaining and describing at B1.
\end{grammarlens}

\subsection*{Your turn}
\begin{itemize}
\raggedright
  \item \emph{Say it:} \emph{brindar}
  \item \emph{Say it:} \emph{brindar}, once more
  \item \emph{Say it:} a sentence of your own with \emph{brindar}
  \item \emph{Recall:} say the Spanish for to marry, then the Spanish for the husband, then say \emph{brindar} again
\end{itemize}

\begin{tcolorbox}[breakable,colback=teal!4,colframe=teal!35!black,title={Before you move on}]
What does brindar mean? (\textquotedblleft{}To toast, to drink to; to offer\textquotedblright{}.) Say it once more.
\end{tcolorbox}
